\documentclass[12pt]{article}
\usepackage{amsmath, amssymb}
\usepackage{booktabs, threeparttable}
\usepackage{graphicx}
\usepackage{placeins}
\usepackage[margin=1in]{geometry}
\usepackage{setspace}
\usepackage[round]{natbib}
\usepackage{hyperref}
\hypersetup{colorlinks=true, citecolor=blue, linkcolor=blue, urlcolor=blue}
\newcommand{\doi}[1]{\href{https://doi.org/#1}{\texttt{doi:#1}}}
\onehalfspacing

\title{\Large\bfseries Reconstruction Without Reallocation:\\
How Did Chilean Manufacturing Capital Adjust\\
After the 2010 Maule Earthquake Shock?}
\author{Sitong Liu\thanks{Undergraduate Honors Program in Economics, New York University.
This is an independent research project. I thank \dots\ All errors are my own.}}
\date{\today}

\begin{document}
\maketitle

\begin{abstract}
\noindent In February 2010, Chile was struck by the Maule earthquake, which measured 8.8 on the
moment magnitude scale and dealt a severe blow to the local manufacturing sector. This paper
treats the Maule earthquake as an exogenous shock to Chile's manufacturing sector to examine how
the sector recovered in the aftermath of the earthquake, whether this recovery spilled over into
the labor market, and whether the recovery had heterogeneous effects on unskilled and skilled
workers. I address this question using a difference-in-differences design across industries and
regions. The results show that, among the variables examined---output, employment, labor
productivity, the labor share of output, proxy measures of profits, and investment---only
investment responded to capital reconstruction, while no significant changes were observed in the
other indicators. However, since the severely affected central urban regions already had higher
investment levels than other areas prior to the disaster, I interpret this evidence as suggestive
rather than conclusive. The paper primarily demonstrates that reconstruction restored the capital
stock but did not alter the scale of production or the allocation of resources.

\medskip
\noindent \textbf{Keywords:} natural disasters, capital-skill complementarity, labor markets,
Chile. \quad \textbf{JEL:} J23, J31, Q54, O14, R11.
\end{abstract}

\thispagestyle{empty}
\clearpage
\setcounter{page}{1}

\section{Introduction}\label{sec:intro}

In the early hours of February 27, 2010, the 8.8-magnitude Maule earthquake struck Chile. It was
the sixth-largest earthquake in recorded history and the second strongest in Chile's history
\citep{barcena2010}. More than 100 aftershocks with magnitudes of 5.0 or higher followed. Over
1.8 million people were affected, infrastructure across the country suffered severe damage, and
economic losses reached \$30 billion, accounting for approximately 18\% of Chile's GDP
\citep{barcena2010}. The three hardest-hit regions, Maule, B\'io-B\'io, and O'Higgins, constitute
the heartland of Chile's manufacturing sector. Together, these regions account for more than a
quarter of Chile's manufacturing output, spanning industries such as food processing, beverages,
steel, paper, shipbuilding, and oil refining, while also contributing approximately half of the
country's agricultural and forestry output \citep{barcena2010}. According to a study using
synthetic control methods, economic activity in the Maule Region declined by approximately 5\%
compared to a counterfactual scenario in the three years following the earthquake
\citep{diaz2024}. Figure~\ref{fig:map} shows the distribution of seismic intensity.

As can be seen, this earthquake was plausibly exogenous to the industrial and labor markets; it
had a massive impact, and the intensity of ground shaking it generated attenuated with distance
from the epicenter and varied with local soil conditions. Therefore, after controlling for fixed
effects, the intensity of shaking is plausibly uncorrelated with a region's pre-existing
manufacturing and labor-market trends. I can thus attempt to answer the following question: How
does a labor market adjust when an economy suddenly loses a large amount of physical capital
overnight due to an exogenous disaster? \citet{krusell2000} proposed the ``capital-skill
complementarity (CSC)'' hypothesis to explain the relationship between equipment capital and the
labor market: equipment capital and skilled labor are complementary in production. I can thus
propose the initial hypothesis that, after an earthquake destroys most of the local physical
capital, the local economy will require more skilled workers to rebuild capital, thereby driving
up wages for skilled workers.

To test whether this theory applies to this earthquake disaster, I aggregate establishment-level
data into industry $\times$ region $\times$ year cells and construct a difference-in-differences
specification with an interaction term between the time indicator variable (post-2010) and
continuous regional MMI values. The industry $\times$ region fixed effects capture the
time-invariant characteristics of each cell, including MMI levels; the industry $\times$ year
fixed effects capture any nationwide industry-specific shocks, such as changes in copper prices,
global demand for processed foods, and the national business cycle. Thus, I compare establishments
across regions within the same industry and year that were affected by the earthquake to varying
degrees.

Reconstruction unfolds over several years rather than instantaneously. The immediate aftermath is
dominated by relief, damage assessment, and financing rather than rebuilding, so the investment
response should appear with a delay, peaking in 2014--2015 rather than at the moment of the quake.
Consistent with a slow recovery, \citet{diaz2024} find that Chilean per-capita GDP remained
roughly 5\% below its counterfactual for years after the earthquake. This motivates the central
question of the paper: during reconstruction, how did the manufacturing labor market adjust,
through capital, through workers, or through wages?

The paper makes three points in response to the above questions.

\textbf{(1) A reconstruction-investment pulse.} Areas that were more severely affected by the
earthquake saw an increase in total investment following the disaster, peaking in 2014 before
declining. For every one-unit increase in a region's MMI, the post-earthquake peak investment in
that region rose by approximately 0.58 log points. This estimate withstood a wild cluster
bootstrap and a series of robustness tests. However, it was found to be sensitive to the faster
pre-earthquake investment growth trend in affected regions.

\textbf{(2) Reconstruction occurred after the earthquake, but the allocation of resources did not
change.} Investment was the only variable to show a significant change post-earthquake. I did not
observe significant changes in other variables such as output, employment, labor productivity, the
labor share of output, or proxy indicators of corporate profits (see Table~\ref{tab:incidence}).
Therefore, this suggests that reconstruction merely replaced capital destroyed by the earthquake
rather than adding new productive capacity, and thus did not alter the allocation structure between
capital and labor. This conclusion is consistent with the fact that Chile's value added remained
below its counterfactual path several years later \citep{diaz2024}. This is the key finding of the
paper and the fundamental reason why the effects of post-earthquake reconstruction did not spill
over into the labor market.

\textbf{(3) The expected ``capital-skill complementarity'' theory has not been confirmed.} If
reconstructed capital and skilled labor were complementary, the post-disaster skill premium should
have risen; however, this pattern did not materialize following the Chilean earthquake. Yet my aim
is not to refute this theory, as the available data and the reduced-form regression I employ do not
solely examine the impact of equipment on the skill premium; they also account for the effects of
variables such as structures, labor, and output---all of which were similarly affected by the
earthquake---on the skill premium. Furthermore, the wage evidence is not sufficiently robust. I
find a weak and imprecise upward trend in unskilled wages following the disaster. A candidate
mechanism is that local reconstruction of destroyed buildings and infrastructure after the disaster
led to increased employment in the construction sector. I therefore used household data to test
this hypothesis but found no significant results.

\textbf{Pre-trends and the identification strategy.} When I regressed the various outcomes on
``post-earthquake $\times$ regional intensity,'' I found that the hardest-hit areas were precisely
Chile's industrial heartland (Santiago, Maule, and B\'io-B\'io), which had already been experiencing
faster investment growth than outlying regions prior to the earthquake. To address this issue, I
recovered the ENIA files with establishment identifiers from before 2008 and extended the
pre-earthquake period back to 2003 to capture regional trends. I then incorporated region-specific
trends estimated solely within the pre-earthquake window. Even so, this study cannot completely
disentangle pre-earthquake investment growth from post-earthquake investment effects, so the
conclusions presented here are tentative.

\textbf{Related literature.} In the economics of natural disasters, a significant body of research
focuses on regional output following major shocks, as well as how populations adapt and recover
\citep{aksoy2024,diaz2024,benguria2022}, while a smaller body of research examines labor market
conditions in affected regions. The study most closely aligned with this paper is
\citet{belasen2009}, who found that hurricanes in Florida increased local income but reduced
employment. Building on these studies, this paper also examines how manufacturing establishments
adjust their total investment, asset mix, and occupational composition following a disaster.

Second, some studies have used capital cost shocks to investigate the extent to which such shocks
are passed on to workers \citep{suarezzidar2016,fuest2018}, while \citet{notowidigdo2020} examined
how changes in local labor demand affect skill groups differently. I adopt a similar logic, but
with fundamental differences from previous studies, which also lead to different results. The
shocks in this paper involve physical destruction, such as damage to factory machinery, buildings,
and infrastructure, rather than changes in tax rates or the cost of capital. Furthermore, I find
that the physical destruction caused by earthquakes is rarely passed on to workers; even when it
is, the pass-through is not significant.

Third, this paper provides an empirical case study of ``capital-skill complementarity.'' The
capital-skill complementarity theory proposed by \citet{krusell2000} emphasizes that capital
equipment increases the relative demand for skilled labor; meanwhile, \citet{lewis2011} argues that
the local labor supply, in turn, shapes an establishment's choice of equipment. The Chilean
earthquake provides an empirical case: in the aftermath of the disaster, local manufacturers
increased investment in equipment to restore the economy, yet the skill premium did not rise.
However, this conclusion is based on the results of a reduced-form regression and is not entirely
robust. Therefore, it is insufficient to refute the validity of Krusell's theory in this specific
case.

Finally, this paper draws on the pre-trends approach of \citet{roth2022} to address the issue that
the hardest-hit regions already experienced faster investment growth prior to the disaster.
Section~\ref{sec:framework} presents the concepts and theories discussed in this paper;
Section~\ref{sec:data} introduces the data; Section~\ref{sec:strategy} outlines the empirical
strategy; Sections~\ref{sec:results}--\ref{sec:mechanism} present the results and mechanisms;
Sections~\ref{sec:migration}--\ref{sec:honest} demonstrate the exclusion of the out-migration
channel and the treatment of parallel trends; Section~\ref{sec:discussion} discusses the effects
and external validity; and Section~\ref{sec:conclusion} concludes the paper.

\begin{figure}[htbp]\centering
  \includegraphics[width=0.5\linewidth]{../output/figs/E2_chile_mmi_map_zoom.png}
  \caption{Ground-shaking intensity (MMI) of the 2010 Maule earthquake by region; the red star
  marks the offshore epicenter. Central regions were severely shaken; the far north and south
  were essentially unaffected.}
  \label{fig:map}
\end{figure}

\FloatBarrier
\section{Conceptual Framework}\label{sec:framework}

According to the capital-skill complementarity theory proposed by \citet{krusell2000}, let $k_s$
denote structures (buildings used for production), $k_e$ denote equipment, $s$ denote skilled
labor, and $u$ denote unskilled labor. The complementarity between equipment and skilled labor is
stronger than that between equipment and unskilled labor:
\begin{equation}
Y = A\,k_s^{\alpha}\left[\mu\,u^{\sigma}+(1-\mu)\left(\lambda\,k_e^{\rho}+(1-\lambda)\,s^{\rho}\right)^{\sigma/\rho}\right]^{(1-\alpha)/\sigma}.
\end{equation}

$\rho$ measures the ease of substitution between equipment and skilled labor; the smaller it is,
the harder it is for the two to substitute for one another. $\sigma$ measures the ease of
substitution between equipment and unskilled labor; the larger it is, the easier it is for the two
to substitute for one another. If $\sigma > \rho$, then equipment is harder to substitute for
skilled labor than for unskilled labor. In other words, machines require highly skilled workers to
operate them, but are less selective when it comes to unskilled workers.

We can therefore conclude that the relative marginal product of skilled labor rises as the stock of
equipment increases; consequently, the skill premium $\omega = w_s/w_u$ will rise as the amount of
equipment allocated to skilled workers increases:
\begin{equation}
\frac{\partial \ln\omega}{\partial \ln k_e}>0 \qquad\text{iff}\qquad \sigma>\rho.
\end{equation}

Equation (2) provides a standard prediction: if equipment stock increases, the wage premium of
skilled labor relative to unskilled labor will increase accordingly. Since plant capital is treated
as a separable Cobb-Douglas factor within Equation (1), it does not carry this implication, which
is consistent with the ``capital-specificity'' noted by Krusell et al.

This baseline formula overlooks a key feature of post-disaster reconstruction. Reconstruction itself
constitutes a form of production activity, involving labor-intensive manual and construction work,
such as demolition, site clearance, and the installation and construction of new facilities and
equipment, which relies heavily on unskilled workers. Therefore, when severely affected regions
rely more heavily on post-disaster facility reconstruction, they have a greater need for unskilled
workers, leading to an outward shift in the demand for unskilled labor, where $\Delta\ln u^{d} > 0$.
However, the supply of unskilled labor is not infinitely elastic; this rise in demand will push
$w_u$ upward, causing the skill premium we obtain to be underestimated, because of the opposite of
the direction of effect in Equation (2). Consequently, the net effect on $\omega$ is theoretically
uncertain:
\begin{equation}
\underbrace{\frac{\partial \ln\omega}{\partial \ln k_e}}_{>0\ (\text{CSC})}\qquad\text{vs.}\qquad
\underbrace{-\,\frac{\partial \ln\omega}{\partial \ln u^{d}}}_{<0\ (\text{reconstruction demand})}.
\end{equation}

This framework raises two empirical questions that this paper seeks to address: H1 (Investment):
Does the shock increase investment in the affected region? H2 (Allocation): Where does this
investment manifest itself---in output, employment, factor allocation, or skill premiums? Equation
(2) provides an expectation that this investment will ultimately manifest itself in rising skill
premiums. However, Equation (3) casts doubt on this expectation. Because the shock simultaneously
alters $k_e$, $k_s$, $u^{d}$, and output, the reduced-form regression performed in this paper
captures the net effect of $k_e$, $k_s$, $u^{d}$, and output as a whole, rather than a pure,
isolated $\partial\ln\omega/\partial\ln k_e$. The data confirm H1, but the result for H2 is that
the investment is not reflected in the actual marginal rate of return nor in factor allocation.
Therefore, I argue that post-earthquake reconstruction was limited to restoring destroyed capital
and infrastructure, without altering the total volume of production or the pattern of allocation
between labor and capital. The prediction of a skill premium was not confirmed, nor was it supported
by existing data for a clean test.

\FloatBarrier
\section{Data}\label{sec:data}

\subsection{Manufacturing microdata}
The Annual National Industrial Survey (\textit{Encuesta Nacional Industrial Anual}, ENIA) is
conducted annually by the Chilean National Institute of Statistics (INE). The survey covers
manufacturing establishments with ten or more employees. I use survey data from 2003 to 2017,
excluding 2012 because the response rate from establishments was unusually low that year. Each
survey includes data reported by establishments on employment, wages and salaries by occupational
category, gross production value (\textit{valor bruto de la producci\'on}, VBP), intermediate
consumption, taxes, and fixed-asset acquisitions and book values by asset category.
Appendix~\ref{app:data} details the data processing procedures.

Three points warrant emphasis. First, public-use files from 2008 onward lack stable establishment
identifiers, so establishment-level data is unavailable; I can only analyze industry-region-year
cells. Additionally, the files I reconstructed for the period prior to 2008 contain stable
identifiers (NUI), which can be used to extend the data to the pre-earthquake period. Second, ENIA
does not directly survey capital stock; therefore, the data do not directly reflect how much capital
was destroyed by the earthquake or how much was replenished. Consequently, I directly use annual
acquisition expenditures as an indicator of capital formation. I categorize these expenditures into
structures (land and buildings), machinery, and a broad equipment measure (machinery, vehicles, and
software). Third, starting in 2013, the industrial classification standard used in the statistics
changed from ISIC Rev.\ 3 to Rev.\ 4. I use the official UN concordance table to reconcile the coding
differences between the two versions, prioritizing the UN standard and classifying industries with
discrepancies based on their weights.

\subsection{Skill, wages, and the limits of the employment data}
To test the theory of ``capital-skill complementarity,'' I must first divide workers into two
categories: skilled workers and unskilled workers. ENIA reports employment figures and wages by
occupational category, so in theory, I could calculate the proportion of skilled workers based on
the number of workers in each category across regions. However, ENIA's classification of these two
categories is somewhat vague and does not provide a clear distinction between unskilled and skilled
workers. Fortunately, the wage category is broken down into monthly wages of technical/professional
workers and monthly wages of unskilled workers, so I use these two wage series to calculate the
skill premium. The skill premium is defined as the log of the ratio of the monthly wages of
technical/professional workers to the monthly wages of unskilled workers.

\subsection{Earthquake intensity}
I use the Modified Mercalli Intensity (MMI) calculated by the USGS ShakeMap \citep{usgs2011} to measure earthquake
intensity. Unlike the magnitude of 8.8, which represents the total energy of the entire earthquake,
the MMI indicates the actual intensity of shaking felt at a specific location. This allows me to
distinguish between severely affected and lightly affected areas based on the intensity of the
shaking. By overlaying the intensity contours from ShakeMap onto administrative district maps and
then applying area-weighted averaging and an equal-area projection, I can calculate the average
intensity for each administrative district. The average MMI across regions ranges from approximately
2.9 (in the northern desert regions) to 6.7 (in Maule and O'Higgins). Areas that were virtually
unaffected by the earthquake are assigned the lower-limit intensity of 2.0 and serve as the control
group. To ensure consistency in regional coding across the data, I use the classification standard
of the 15 major regions. Therefore, the \~Nuble region is still grouped under the B\'io-B\'io region.
I treat continuous intensity as the main explanatory (treatment) variable and use a binary indicator
variable for high intensity (MMI $\geq 6$) as a robustness test.

\subsection{Sample and descriptive statistics}
I aggregate the data into industry $\times$ region $\times$ year cells, calculate the weighted sum
of aggregate economic indicators such as investment and employment, and compute the weighted average
of wages. I select food products (ISIC 10), wood products (ISIC 16), and basic metals (ISIC 24) as
the sample industries because these three sectors are widely distributed and have comparable levels
of development across both low- and high-intensity regions, ensuring that the comparison between
high- and low-intensity regions is not confounded by differences in industry composition. I report
the results of the robustness check using the entire manufacturing sample.
Table~\ref{tab:summary} presents the conditions in severely and lightly affected areas prior to the
earthquake. The results indicate that the severely affected areas were substantially ahead of the
lightly affected areas in terms of investment, employment, and output. This suggests that the most
severely affected central regions were already the industrial heartland of Chile. Therefore, I
suspect that the two groups of regions were likely already following different growth trajectories
prior to the earthquake. This is why I estimate pre-earthquake trends and project them forward in
Section~\ref{sec:strategy}, and test their robustness in Section~\ref{sec:honest}.

\input{tables/tab_summary}

\FloatBarrier
\section{Empirical Strategy}\label{sec:strategy}

\subsection{Specification}
The baseline difference-in-differences specification is
\begin{equation}
y_{jrt} = \beta\,\big(\text{Post}_t \times \text{MMI}_r\big) + \gamma_{jr} + \delta_{jt} + \varepsilon_{jrt},
\end{equation}
where $j$ indexes 2-digit industry, $r$ region, and $t$ year; $\text{Post}_t = \mathbf{1}[t \geq 2010]$;
$\text{MMI}_r$ is the region's area-weighted intensity; $\gamma_{jr}$ are industry $\times$ region
fixed effects; and $\delta_{jt}$ are industry $\times$ year fixed effects. The industry $\times$
region effects absorb the main effect of MMI and any time-invariant cell heterogeneity; the industry
$\times$ year effects absorb any shock common to an industry nationwide each year. $\beta$ is the
coefficient of interest: it captures the differential change in the outcome after 2010 associated
with a one-unit-higher MMI in a region---that is, how much more the outcome moved after the earthquake
in regions shaken one unit harder. I replace $\text{MMI}_r$ with a binary variable (indicating whether
an area is severely affected) as a robustness check, and weight the cells by size, with standard errors
clustered by region.

\subsection{Identification, pre-trends, and detrending}
However, Equation (4) is constrained by the parallel-trends assumption of the difference-in-differences
method, which holds that if the earthquake had never occurred, development in regions with high MMI and
low MMI would have followed parallel trends within each industry. Yet, with 2009 as the base year and
2003--2009 as the pre-earthquake phase, this parallel-trends assumption is fragile in this case, as
regions with high MMI already had higher levels of investment, employment, and output prior to the
earthquake. Equation (5) is designed to test this assumption:
\begin{equation}
y_{jrt} = \sum_{k\neq 2009} \beta_k\,\big(\mathbf{1}[t=k] \times \text{MMI}_r\big) + \gamma_{jr} + \delta_{jt} + \varepsilon_{jrt}.
\end{equation}
To reduce the influence of these pre-existing trends, I estimate each region's trend over the
pre-earthquake period (2003--2009), project it forward to construct a counterfactual path, and subtract
this counterfactual from the actual data. Then I use the resulting residuals to estimate Equation (5).
I also use the 2014 FUT tax reform as a placebo shock for a falsification test.

\subsection{Inference with few clusters}
Because there are only 15 regions, cluster-robust standard errors are downward biased and the asymptotic
$t$-test over-rejects. For the headline coefficients I therefore also report restricted wild cluster
bootstrap $p$-values (Rademacher weights, $B = 999$, the null imposed), following \citet{cameron2008}.

\FloatBarrier
\section{Results}\label{sec:results}

\subsection{The reconstruction-investment pulse}
Table~\ref{tab:main_did} reports the baseline coefficient from (4) for each outcome. Panel A, without
region trends, shows large and significant effects on investment, gross output, and employment. Panel
B adds region-specific linear trends: the investment coefficient falls from 0.52 to 0.17 and loses
significance, and the output and employment effects disappear. Most of Panel A is the differential
secular growth of the central regions, not the earthquake, and this result is exactly the concern
raised by the long pre-period.

\input{tables/tab_main_did}

However, for a transient phenomenon, a single post-earthquake coefficient is not an appropriate
measure, as it reduces the results to an average, preventing me from seeing how the results change over
time. Therefore, I conducted an event study organized by year. Table~\ref{tab:event_study} and
Figure~\ref{fig:eventstudy} present the results of the event study. Even during the long pre-earthquake
period, the investment gap in high-MMI regions remained positive relative to the base year; therefore,
the event study remains inconclusive and cannot be shown to strictly follow a parallel pre-event trend.
However, it shows that during the reconstruction period, the divergence in investment gaps between
severely and lightly affected areas became very large: the investment gap between high and low MMI
areas rose to 0.58 in 2014 ($p < 0.01$) and to 0.60 in 2015 ($p < 0.05$), before falling back to near
zero in 2016--2017.

These changes track the institutional structure of the reconstruction. Chilean authorities adopted a
four-year reconstruction plan, financed by temporarily raising corporate tax rates, increasing mining
royalties, drawing on the sovereign wealth fund, and issuing new bonds \citep{govchile2010}. The
Ministry of Housing and Urban Development (MINVU) committed to rebuilding or repairing 220,000 housing
units over four years, but noted that the affected population and housing stock were geographically
dispersed across rural areas and that land-ownership and technical constraints made rapid execution
difficult. The immediate post-disaster priorities were therefore relief, assessment, subsidies, and
approvals; manufacturing establishments recorded the acquisition and replacement of structures,
machinery, vehicles, and software only once projects reached the implementation phase. The 2014--2015
investment peak thus reflects the execution of reconstruction rather than an immediate response to the
physical shock.

The magnitude of the investment effect in 2014 depends on the method of trend removal: the effect is
0.71 with no detrending, 0.40 with detrending applied to the entire sample, and I tend to use the
figure of 0.58, which lies between the two and passed the bootstrap test for wild clustering
($p=0.027$).

Table~\ref{tab:robustness} shows that the 2014 coefficient is robust across alternative
specifications---adding industry $\times$ region trends, restricting to a tight 2007--2017 window,
dropping the Santiago Metropolitan Region, and using the full manufacturing sample, where it ranges
from 0.56 to 0.88 and remains significant. Under a binary high-intensity (MMI $\geq 6$) treatment it is
larger still (1.87) but also significant. This supports H1.

\input{tables/tab_event_study}
\input{tables/tab_robustness}

\begin{figure}[htbp]\centering
  \includegraphics[width=0.9\linewidth]{../output/figs/E5_event_study_pretrend.png}
  \caption{Investment event study with region trends estimated on the pre-period (2003--2009) only
  and projected forward. The 2014--2015 spike is substantially larger than the remaining pre-period
  gaps and reverts thereafter.}
  \label{fig:eventstudy}
\end{figure}

\FloatBarrier
\subsection{Reconstruction without reallocation}
Table~\ref{tab:incidence} addresses the question posed by H2: Where does the increase in investment
manifest itself? The answer is: nowhere. I apply the same pre-period-projected detrending and regression
to output, employment, labor productivity, the labor share of output, and the corporate-income-tax
profit proxy. The results are all statistically indistinguishable from zero. In other words,
reconstruction restored the destroyed capital but did not expand production, alter the level of
employment or labor productivity, or change the distribution of income between labor and capital. I
interpret this as ``replacement'': establishments rebuilt the capital destroyed by the earthquake,
restoring it to its pre-earthquake level rather than expanding beyond the original level. This explains
why the real and distributional margins remained unchanged, and why (consistent with \citealp{diaz2024})
regional value added consistently fell short of the counterfactual path. The only exception is a slight
increase in the wage bill ($+0.14$, $p < 0.05$). I analyze the possible mechanisms behind this wage
increase in Sections~\ref{sec:mechanism} and~\ref{sec:migration}. These include positive spillover
effects from the local construction industry and the out-migration of unskilled workers.

\input{tables/tab_incidence}

\FloatBarrier
\subsection{The skill premium}
Based on the capital-skill complementarity principle, one might expect that the reconstruction of
infrastructure following an earthquake would be accompanied by a rise in the skill premium. However,
Table~\ref{tab:labor} does not support this prediction: the premium effect is negative, averaging
$-0.19$ log points per unit of MMI. The overall wage increase shown in Table~\ref{tab:incidence} stems
primarily from a rise in unskilled wages, averaging $+0.26$ and peaking at $+0.52$ in 2014, while
skilled wages remained stable. This approach has a limitation: it remains a reduced-form regression that
captures the combined impact of equipment, structures, labor, and output on the skill premium in the
wake of the earthquake, rather than the independent effect of any single factor.

\input{tables/tab_labor}

\FloatBarrier
\section{Replacement, and the Missing Wage Mechanism}\label{sec:mechanism}

This section primarily clarifies two points: First, the investment pulse following an earthquake
represents a replacement of destroyed capital rather than further growth from the original level.
Second, my analysis using household data reveals that the conclusion regarding an upward shift in
unskilled wages is fragile. Taken together, these two points indicate that investment is linked to the
demand for post-disaster reconstruction, rather than to a capital-driven shift in labor demand.

\textbf{Investment represents a replacement.} Table~\ref{tab:mechanism} examines the asset classes to
which the increase in investment in 2014 was allocated: structures accounted for the largest share
($+0.82$), while machinery ($+0.49$) and broad equipment ($+0.52$) also saw significant increases (see
Figure~\ref{fig:asset}). Panel B shows that the equipment coefficient was 0.34 higher than its
pre-period trend during both the 2010--2014 and 2015--2017 periods, while the structures coefficient
was lower in the latter period. Therefore, this table serves as key evidence supporting the paper's
argument: reconstruction covered structures and equipment, but did not affect the distribution of
output, employment, or skills, primarily because reconstruction merely restored the capital base rather
than expanding or developing it beyond its original level.

\input{tables/tab_mechanism}

\begin{figure}[htbp]\centering
  \includegraphics[width=0.9\linewidth]{../output/figs/E8_asset_decomposition.png}
  \caption{Reconstruction investment pulse by asset class. Both structures and equipment rise
  significantly; the null on skill is therefore not due to an absence of equipment investment.}
  \label{fig:asset}
\end{figure}

\FloatBarrier
\textbf{The rise in unskilled wages did not originate in the construction industry.} As shown in
Figure~\ref{fig:wagedecomp}, the only labor-side variable that changed was unskilled wages, which gradually
increased and peaked alongside investment in 2014, while skilled wages remained stable. One plausible
mechanism is the spillover effect from the local labor market: following the disaster, there was a
greater need for unskilled labor to rebuild infrastructure and factories. Consequently, the
construction boom in the hardest-hit areas would have stimulated wages for local unskilled workers,
forcing manufacturing plants to follow suit with wage increases. I directly tested this hypothesis
using household survey data and found that, in the more severely affected regions, post-earthquake
employment in the construction sector did not increase ($+0.02$, $p=0.36$), nor did construction wages
rise. On the contrary, they declined ($-0.04$, $p=0.09$). The event study also revealed no break at the
time of the earthquake. At the same time, I found that the rise in unskilled wages stemmed primarily
from the food industry, rather than the wood products or basic metals industries, the two sectors
closely related to construction. The conclusions are presented in Table~\ref{tab:wage_mechanism_checks}.

\input{tables/tab_wage_mechanism_checks}

\begin{figure}[htbp]\centering
  \includegraphics[width=0.9\linewidth]{../output/figs/E9_wage_decomposition.png}
  \caption{Decomposing the skill-premium decline: the unskilled wage rises and peaks in 2014, in
  lockstep with the investment pulse, while the skilled wage is flat.}
  \label{fig:wagedecomp}
\end{figure}

\FloatBarrier
There are three potential limitations to note. First, the classification criteria for skilled and
unskilled workers in the original data are relatively crude. Second, the skill premium reflects both
labor demand and labor supply: when demand for unskilled labor rises, wages for unskilled workers
increase, and, all else being equal, the skill premium declines. When the supply of unskilled labor
increases (i.e., if more unskilled workers enter the disaster-affected area), wages for unskilled
workers fall, and the skill premium rises. Conversely, if the earthquake causes workers to become
displaced and leave the severely affected areas, the net effect on the skill premium is ambiguous. The
key lies in the post-earthquake mobility of these two categories of workers---specifically, whether they
left the severely affected areas. I discuss this point in Section~\ref{sec:migration}. Third, this paper
draws inferences based solely on data from 15 regions, and this is the primary reason I employed the
bootstrapping method.

\FloatBarrier
\section{The Out-Migration Channel}\label{sec:migration}

There is another possible mechanism to explain why unskilled wages rise: the out-migration effect. If
unskilled workers leave the hardest-hit areas, the local supply of unskilled labor contracts, causing
unskilled wages to rise. However, this mechanism is unrelated to reconstruction. The out-migration
effect and the reconstruction mechanism mentioned earlier appear to have the same directional impact on
wages, while their difference lies in the number of workers: if the rise in unskilled wages is primarily
due to reconstruction, the number of unskilled workers will increase. If it is caused by the
out-migration effect, the number will decrease.

ENIA cannot address this issue because it stopped collecting statistics broken down by skill level after
2007. Therefore, I use another dataset here: the National Socioeconomic Characterization Survey (CASEN).
CASEN records region, educational attainment, and population expansion weights, allowing me to calculate
the stock of skilled and unskilled workers in each region before and after the earthquake. CASEN survey
data is only available for 2003, 2006, 2009, 2011, 2013, 2015, and 2017, with a lower frequency of data
collection than ENIA. Therefore, I primarily use it for robustness checks. I adopt the same
specifications as in the main analysis ($\text{Post} = \mathbf{1}[t\geq 2010]$, where the 2009 wave is
classified as pre-earthquake) and divide regions into the 13 major regions as defined prior to 2007.
Skill levels are still categorized by educational attainment: unskilled workers have a basic education
or lower, while skilled workers have a higher education background.

Table~\ref{tab:casen} and Figure~\ref{fig:casen} present the results, which support the demand-side
interpretation. The working-age unskilled population in severely affected areas did not decline after
the earthquake: the $\text{Post} \times \text{MMI}$ coefficient was $+0.002$, which is statistically
indistinguishable from zero. Between 2003 and 2017, the unskilled population in areas with high MMI
maintained a small positive difference, and no breakpoint occurred at the time of the shock. The
coefficient for 2006 in the event study was positive, but the combined test of the pre-event trend
(i.e., that the coefficients for both 2003 and 2006 were equal to zero) was not rejected ($F = 2.60$,
$p = 0.12$). To rule out the out-migration confound, I used the same HonestDiD test as in the main
results. The results show that when $\bar M = 1$, the robust confidence interval for the relative
magnitude of the average post-earthquake effect is $[-0.07, 0.12]$. Therefore, even if I allow for a
parallel-trend violation after the earthquake that is as large as the worst-case scenario before the
earthquake, the unskilled population in the more severely affected areas would decline by no more than
approximately 7 percent per unit of MMI. This figure is far insufficient to account for the level of
unskilled wage increases mentioned in Section~\ref{sec:mechanism}. At the same time, the data in
Table~\ref{tab:casen} indicate that the unskilled labor force actually increased slightly after the
earthquake, and its share of the total population remained unchanged. Therefore, an out-migration effect
can be ruled out. Furthermore, the slight decline in the total population ($-0.009$ per unit of MMI) is
attributable not to the unskilled labor force but to the skilled labor force. At the same time, wages for
the skilled labor force did not change significantly despite the decline in the size of this group. The
mechanisms underlying this phenomenon are beyond the scope of this paper and may thus serve as a topic
for future discussion and research. This conclusion still has some data-related limitations: the CASEN
survey is conducted infrequently, and the small number of clusters (13 regions in the harmonized CASEN
scheme) may limit the precision of the estimates.

\input{tables/tab_casen}

\begin{figure}[htbp]\centering
  \includegraphics[width=0.75\linewidth]{../output/figs/E12_casen_migration.png}
  \caption{CASEN: regional population by skill versus earthquake exposure. A fall in the unskilled
  population in harder-hit regions would indicate out-migration; the coefficient is instead near zero.}
  \label{fig:casen}
\end{figure}

\FloatBarrier
\section{Sensitivity to Parallel Trends}\label{sec:honest}

The identification in this paper relies on the parallel-trend assumption, namely, that after controlling
for region-specific trends, the growth trends of various variables (investment, output, employment,
wages/skill premiums, etc.) in regions with high MMI and low MMI are parallel. If the pre-shock
coefficients in the regression model are approximately equal to zero and statistically insignificant,
then the pre-shock parallel-trend assumption holds. However, I found that not all pre-earthquake
coefficients were equal to zero. To address this issue, I employed the sensitivity analysis method
proposed by \citet{rambachanroth2023} to assess the core empirical conclusion of this paper, namely, to
what extent the investment spike in 2014, resulting from post-earthquake reconstruction, depends on the
assumption of prevailing trends. This method moves beyond the original binary classification of
``parallel'' versus ``non-parallel'' to determine the extent to which the pre-earthquake trend can
deviate from parallelism while the 2014 effect remains statistically significant. I applied this
approach to the raw event-study estimates, summarizing the results in Table~\ref{tab:honestdid} and
Figure~\ref{fig:honest}.

There are two natural ways to impose constraints here. The smoothness constraint $\Delta^{SD}(M)$ limits
the change in the slope of the divergence trend in each period to no more than $M$, meaning that the
trend line is not allowed to turn sharply and can only change gradually. When $M=0$, the trend is forced
to be a straight line, which is entirely consistent with the linear detrending process I described
earlier. The relative magnitude constraint $\Delta^{RM}(\bar M)$, on the other hand, uses the actual
deviations observed before the earthquake as a baseline and limits the magnitude of deviations from the
post-earthquake parallel trend to within $\bar M$ times the maximum pre-earthquake deviation. The
advantage of this approach is that it is not affected by the uneven distribution of years caused by the
exclusion of the 2012 wave. These two constraints operate from the perspectives of ``how sharply the
trend can turn'' and ``how large the post-earthquake deviation is relative to the pre-earthquake value,''
respectively. I apply both simultaneously to ensure robustness.

Under the linear-trend benchmark $\Delta^{SD}(0)$, the robust confidence sets for the investment
coefficient and the unskilled wage coefficient in 2014 both exclude zero ($[0.28, 0.75]$ and
$[0.19, 0.80]$, respectively), consistent with the estimates preferred in this paper. However, these
conclusions are fragile in the face of larger deviations: the breakdown values for both results are
approximately $\bar M \approx 0.5$, which is lower than the commonly used benchmark of $\bar M = 1$
(under which post-earthquake violations as large as the most severe pre-earthquake violations are
permitted, at which point both confidence sets include zero).

I also used another, more refined estimation method to assess robustness. Since locally smoothed trends
vary approximately linearly between adjacent years, the difference $\beta_{2014} - \tfrac12(\beta_{2013}
+ \beta_{2015})$ can account for such a trend, thereby isolating any ``spikes'' in a single year.
However, this does not help in this case (Table~\ref{tab:honestdid}, last row): the difference itself is
not significant because the effect is a multi-year upward shift between 2013 and 2015, rather than a spike
in a single year; thus, subtracting 2014 from its adjacent years removes the signal along with the trend.
Combined with standard linear and quadratic trend specifications (Table~\ref{tab:main_did}, Panels B--C,
under which the investment effect barely survives), I believe the data support the following conclusion:
as long as deviations from the parallel trend are mild and smooth, the findings regarding the
reconstruction of the investment pulse and the constancy of unskilled wages are solid. However, once
deviations become too large, these findings become fragile.

\input{tables/tab_honestdid}

\begin{figure}[htbp]\centering
  \includegraphics[width=0.95\linewidth]{../output/figs/E11_honestdid.png}
  \caption{HonestDiD sensitivity \citep{rambachanroth2023} of the 2014 coefficients to
  relative-magnitude violations of parallel trends. The robust 95\% confidence set widens with
  $\bar M$; it excludes zero only for $\bar M$ below about 0.5.}
  \label{fig:honest}
\end{figure}

\FloatBarrier
\section{Discussion}\label{sec:discussion}

\textbf{Post-earthquake reconstruction focused primarily on replacing the damaged capital rather than
rebuilding to the original level of development.} Investment increased to rebuild and replace capital,
but there were no changes in output, employment, the allocation of key factors, or wages in any robust
sense. Under these circumstances, the marginal returns to capital and labor returned to pre-earthquake
levels, so neither the prices nor the quantities of these two factors changed significantly. At the same
time, I found that when simulating a counterfactual path in which the earthquake shock did not occur, the
actual level of development remains below that of the counterfactual path. This is consistent with the
findings of \citet{diaz2024}, namely that investment offsets the damage rather than building upon the
original foundation. Furthermore, because investment is not aimed at development, its impact on the labor
force and income distribution is minimal.

\textbf{Alternative explanations.} Theoretically, several other mechanisms could also explain the findings
of this paper. (a)~\emph{Nominal and real wages.} The data used in this paper are based on nominal wages,
so inflation during the reconstruction period may have led me to overestimate changes in unskilled or
skilled wages. However, the skill premium differences out common inflation, so the skill-premium
conclusion is not affected by inflation. (b)~\emph{Insurance and reconstruction transfer payments.} Public
reconstruction spending and insurance payouts boosted consumption and purchasing power in severely
affected areas, leading to increased demand for local goods and services. This may account for part of
the change in local demand, but it does not invalidate my regression results. (c)~\emph{Survivorship
bias.} If larger establishments were more likely to survive the earthquake, while smaller ones were more
likely to disappear and thus not be included in post-earthquake statistics, then the post-earthquake data
would be skewed toward surviving establishments. However, Table~\ref{tab:main_did} shows that the change
in the number of establishments is zero and not statistically significant. (d)~\emph{Output demand for
construction materials.} Reconstruction increases the construction industry's demand for building
materials such as lumber, cement, and metals. However, wage increases were concentrated in the food
industry, the sector least associated with reconstruction. The specific reason for wage increases in the
food industry is not the focus of this paper but may serve as a new research idea for the future.

\textbf{2014 tax reform.} In 2014, the Chilean government enacted a new corporate tax policy (Law 20,780),
which could have led to front-loading---that is, companies rushing to make investments before the tax
reform officially took effect, resulting in a peak in investment in 2014. Moreover, such front-loading was
more likely to occur in disaster-stricken areas where businesses were already more concentrated. However,
the paper has ruled out this possibility on two grounds. First, the assets that establishments would
front-load are those that can be purchased in advance, such as machinery. Therefore, if a 2014 investment
peak were indeed caused by anticipatory front-loading, that peak should have been dominated by equipment.
However, the findings in Table~\ref{tab:mechanism} show that the investment peak was dominated by
structures, which establishments cannot purchase in advance. Second, no further spikes should have
occurred after the tax reform was implemented, so the investment peak should have subsided in 2014.
However, my data show that the investment peak persisted into 2015. Therefore, I believe that the 2014 tax
reform was not the primary cause of the investment peak and does not affect the findings of this paper.

\textbf{Spillover effects.} In this paper, the severely affected areas with high MMI are treated as the
treatment group, while the lightly affected areas with low MMI are treated as the control group. However,
the lightly affected areas were not entirely unaffected. Post-earthquake reconstruction requires large
quantities of cement, lumber, steel, and workers, all of which are likely to be transferred from the
lightly affected areas to the severely affected areas. In other words, because the lightly affected areas
are supplying materials to the severely affected areas, manufacturing demand rises, pulling up output and
investment, which narrows the gap between the treatment group and the control group, and causes $\beta$ to
underestimate the true effect. At the same time, the transfer of labor from the lightly affected areas to
the severely affected areas reduces employment in the former, widening the gap between the control group
and the treatment group, and that causes $\beta$ to overestimate the true effect.

I test for the spillover directly with a spatial-lag-of-treatment (SLX) specification, adding a second
treatment term, $\text{Post} \times \text{MMI}$ for a region's immediate geographic neighbors, to the
incidence regressions. Table~\ref{tab:spillover} reports the result, and the result is that the spillover
cannot be signed with this design. Because Chile is long from north to south but very narrow from east to
west, its geography is almost linear, with the various regions essentially lined up end-to-end.
Therefore, a region's own shaking and its neighbors' shaking are nearly collinear ($r \approx 0.98$), so
the own- and neighbor-exposure terms are not separately identified. The investment neighbor coefficient is
0.51 with a standard error above one, $p = 0.64$, and the output own-coefficient inflates from a precise
null of 0.11 to 1.06 while a mechanically ``significant'' neighbor term of $-0.97$ appears with the
opposite sign. This means the standard errors explode and the augmented own-coefficients become unstable.
Therefore, the sign of the spillover is ambiguous.

\input{tables/tab_spillover}

\textbf{Limitations.} (i) Since 2008, the establishment identifiers in the ENIA data have changed
annually, making it impossible for me to identify the same establishment across different years.
Consequently, I am unable to construct an establishment panel and can only aggregate the data into cells
organized by ``industry $\times$ region $\times$ year.'' Therefore, based on the available data, I cannot
distinguish whether the effects I identify stem from internal adjustments within a specific establishment
or from overall changes in the industry average. (ii) ENIA applies multiplicative confidentiality noise,
which increases the variance. To address this, I performed a bootstrap analysis. (iii) Wages for unskilled
workers rose during the post-earthquake reconstruction period, while wages for skilled workers remained
largely unchanged. The specific mechanism behind this phenomenon has not yet been elucidated. Only two
mechanisms have been ruled out: the local construction industry development mechanism mentioned in
Section~\ref{sec:mechanism} and the out-migration effect of unskilled labor mentioned in
Section~\ref{sec:migration}. Further possible mechanisms remain to be explored in future research.
(iv) Regardless of the type of publicly available data, none of it covers the working hours of individual
workers. Therefore, the analysis of the labor force in this paper is limited to the single dimension of
``number of workers.''

\FloatBarrier
\section{Conclusion}\label{sec:conclusion}

I used the 2010 Maule earthquake in Chile as an exogenous shock to examine its impact on Chilean
manufacturing and how the sector adjusted during post-disaster reconstruction. The results show that
regions most severely affected by the disaster experienced a temporary spike in investment, peaking in
2014. At the same time, there were no significant changes in output, employment, labor productivity, the
labor share of output, or the profit proxy. The rise in wages for unskilled workers was the only labor
market change observed, but the mechanism behind it remains unclear. Furthermore, this paper points out
that post-earthquake reconstruction merely restored the capital destroyed by the earthquake, rather than
building upon the existing foundation. Therefore, the post-earthquake investment peak represents
``reconstruction'' rather than ``labor reallocation.''

\clearpage
\begin{thebibliography}{99}
\bibitem[Aksoy et al.(2024)]{aksoy2024} Aksoy, C.~G., Chupilkin, M., Koczan, Z., \& Plekhanov, A.
(2024). Unearthing the economic and social consequences of earthquakes. \textit{IZA Discussion
Paper No.\ 17108}.
\bibitem[B\'arcena et al.(2010)]{barcena2010} B\'arcena, A., Prado, A., L\'opez, L., \& Samaniego, J. (2010). \textit{The Chilean
Earthquake of 27 February 2010: An Overview}. Santiago: Economic Commission for Latin America and
the Caribbean (ECLAC).
\bibitem[Belasen \& Polachek(2009)]{belasen2009} Belasen, A.~R., \& Polachek, S.~W. (2009). How
disasters affect local labor markets: The effects of hurricanes in Florida. \textit{Journal of
Human Resources}, 44(1), 251--276. \doi{10.3368/jhr.44.1.251}
\bibitem[Benguria(2022)]{benguria2022} Benguria Garibaldi, S. (2022). Shaky growth: Chile's
earthquake and its effect on GDP. Master's thesis, Stockholm School of Economics.
\bibitem[Cameron et al.(2008)]{cameron2008} Cameron, A.~C., Gelbach, J.~B., \& Miller, D.~L.
(2008). Bootstrap-based improvements for inference with clustered errors. \textit{Review of
Economics and Statistics}, 90(3), 414--427. \doi{10.1162/rest.90.3.414}
\bibitem[D\'iaz et al.(2024)]{diaz2024} D\'iaz, D., Paniagua, P., \& Larroulet, C. (2024).
Earthquakes and the wealth of nations: The cases of Chile and New Zealand. \textit{Working
paper}, arXiv:2405.12041.
\bibitem[Fuest et al.(2018)]{fuest2018} Fuest, C., Peichl, A., \& Siegloch, S. (2018). Do
higher corporate taxes reduce wages? \textit{American Economic Review}, 108(2), 393--418.
\doi{10.1257/aer.20130570}
\bibitem[Government of Chile(2010)]{govchile2010} Government of Chile. (2010). \textit{Plan de
Reconstrucci\'on Chile Unido Reconstruye Mejor} [A United Chile Rebuilds Better: National Reconstruction Plan]. Santiago: Gobierno de Chile.
\bibitem[Krusell et al.(2000)]{krusell2000} Krusell, P., Ohanian, L.~E., R\'ios-Rull, J.-V.,
\& Violante, G.~L. (2000). Capital-skill complementarity and inequality.
\textit{Econometrica}, 68(5), 1029--1053. \doi{10.1111/1468-0262.00150}
\bibitem[Lewis(2011)]{lewis2011} Lewis, E. (2011). Immigration, skill mix, and capital-skill
complementarity. \textit{Quarterly Journal of Economics}, 126(2), 1029--1069.
\doi{10.1093/qje/qjr011}
\bibitem[Notowidigdo(2020)]{notowidigdo2020} Notowidigdo, M.~J. (2020). The incidence of local
labor demand shocks. \textit{Journal of Labor Economics}, 38(3), 687--725.
\doi{10.1086/706048}
\bibitem[Rambachan \& Roth(2023)]{rambachanroth2023} Rambachan, A., \& Roth, J. (2023). A more
credible approach to parallel trends. \textit{Review of Economic Studies}, 90(5), 2555--2591.
\doi{10.1093/restud/rdad018}
\bibitem[Roth(2022)]{roth2022} Roth, J. (2022). Pretest with caution: Event-study estimates
after testing for parallel trends. \textit{American Economic Review: Insights}, 4(3), 305--322.
\doi{10.1257/aeri.20210236}
\bibitem[Su\'arez Serrato \& Zidar(2016)]{suarezzidar2016} Su\'arez Serrato, J.~C., \& Zidar,
O. (2016). Who benefits from state corporate tax cuts? \textit{American Economic Review},
106(9), 2582--2624. \doi{10.1257/aer.20141702}
\bibitem[U.S. Geological Survey(2010)]{usgs2011} U.S. Geological Survey. (2010). \textit{M 8.8 ---
offshore Bio-Bio, Chile (27 February 2010): Earthquake summary and ShakeMap} (Event ID us2010tfan). Reston, VA: USGS. Retrieved from \url{https://earthquake.usgs.gov/earthquakes/eventpage/us2010tfan}.
\end{thebibliography}

\clearpage
\FloatBarrier
\appendix
\section{Data construction}\label{app:data}

\subsection{Raw ENIA files and sample window}
This paper is based on data from the Annual National Industrial Survey (\textit{Encuesta Nacional
Industrial Anual}, ENIA) collected by the Chilean National Institute of Statistics (INE). The raw data are
collected annually and are divided into three parts: Part 1 records the establishment identifier and
classification information (region, industry, and size); Part 2 records the economic accounts, namely
gross output, fixed-asset acquisitions (investment) by asset class, the wage bill, and taxes; and Part 3
records the detailed labor variables, namely total employment and wages by occupational category. The
main outcome variables in this paper come from Part 2 (investment, gross output, the wage bill, and the
corporate income tax used as a profit proxy); Part 3 provides employment together with the skilled and
unskilled wages behind the skill premium; and Part 1 provides the region and industry codes used to define
the cells.

I first convert each annual database into a single wide-format establishment-level file. The specific
processing steps are: standardizing the establishment identifier, removing duplicate IDs, decoding numeric
fields that were exported as byte strings in 2009, and merging the three parts through an outer join on the
establishment identifier. The working sample covers the years 2003--2017, excluding 2012 because the
establishment response rate to the questionnaire was markedly lower than in other years.

The data used in this paper span two different file regimes. Files prior to 2008 use stable establishment
identifiers and follow the ISIC Rev.\ 3 industry classification. The publicly available FUSION files from
2008 onward use annual identifiers that are not consistent across years; I therefore do not construct an
establishment-level panel for the period after 2007. Instead, all regressions use repeated cross-sections
composed of ``2-digit industry $\times$ region $\times$ year'' cells.

\subsection{Industry harmonization}
The industrial classification standard is not the same across all years. The files I recovered from before
2013 are coded in ISIC Rev.\ 3 (or Rev.\ 3.1), while the later files use ISIC Rev.\ 4. I therefore map the
pre-2013 codes to Rev.\ 4 using the official United Nations concordance table. When one Rev.\ 3 code
corresponds to more than one Rev.\ 4 code, I split the establishment across those Rev.\ 4 industries using
activity-share weights, so that no establishment is double-counted. All of the analysis is then carried out
at the 2-digit Rev.\ 4 level. As in Section~\ref{sec:data}, I restrict the main sample to food products
(10), wood products (16), and basic metals (24), because these three industries appear in both high- and
low-intensity regions and therefore support the within-industry, across-region comparison. I use the full
2-digit manufacturing sample as a robustness check.

\subsection{Cell aggregation and outcome construction}
I aggregate the establishment-level data into cells defined by 2-digit industry, region, and year. Let $i$
index establishments in industry $j$, region $r$, and year $t$, and let $w_i$ be the activity-share weight
from the industry harmonization. For total outcomes, such as investment, employment, output, and the
effective number of establishments, I take the weighted cell sum:
\[ Y^{\mathrm{sum}}_{jrt} = \sum_{i \in (j,r,t)} w_i\, y_i. \]
For wage outcomes, I instead take the weighted mean over establishments with positive, non-missing wages:
\[ Y^{\mathrm{mean}}_{jrt} = \frac{\sum_{i \in (j,r,t)} w_i\, y_i}{\sum_{i \in (j,r,t)} w_i}. \]
I then use the log of the positive cell values as the regression outcomes, except for the skill premium,
which I define as the log skilled wage minus the log unskilled wage.

Table~\ref{tab:varconstruction} lists the main variables and their source fields. The most important
construction issue is investment. Through 2014, purchases of new capital goods are reported in the
\texttt{CBN*} fields; from the 2015 form redesign onward, the same items are reported under \texttt{CB*}
fields. I treat this as a field rename rather than a stop in reporting, and I combine the two series into
one. Any national rescaling or form-level break in 2015 is common to an industry-year and is absorbed by
the industry $\times$ year fixed effects.

\begin{table}[htbp]\centering
\caption{Variable construction and source fields}
\label{tab:varconstruction}
\begin{threeparttable}
\small\setlength{\tabcolsep}{4pt}
\begin{tabular}{p{0.21\linewidth}p{0.42\linewidth}p{0.27\linewidth}}
\toprule
Concept & ENIA source fields & Cell construction \\
\midrule
Effective establishments & Harmonized establishment weights & Weighted sum of $w_i$ \\
Gross investment & \texttt{CBNTER}, \texttt{CBNEDI}, \texttt{CBNMAQ}, \texttt{CBNVEH},
\texttt{CBNMUE}, \texttt{CBNSOF}; from 2015: \texttt{CBTER}, \texttt{CBEDI}, \texttt{CBMAQ},
\texttt{CBVEH}, \texttt{CBSOF} & Weighted cell sum of capital-good purchases \\
Structures & Land and buildings fields, mainly \texttt{CBNTER}, \texttt{CBNEDI} (old schema) and
\texttt{CBTER}, \texttt{CBEDI} (new schema) & Weighted cell sum \\
Machinery & \texttt{CBNMAQ} and \texttt{CBMAQ} & Weighted cell sum \\
Equipment & Machinery, vehicles, and software fields & Weighted cell sum \\
Employment & \texttt{EMPTOT}; for years without \texttt{EMPTOT}, fallback
\texttt{TOTHOM}+\texttt{TOTMUJ} & Weighted cell sum \\
Gross output & \texttt{VBP}; \texttt{VBPb}/\texttt{VBPB} renamed to \texttt{VBP} where needed &
Weighted cell sum \\
Production wage & \texttt{REGPRDI} & Weighted positive cell mean \\
Skilled wage & \texttt{REGTESP} (technicians/specialists) & Weighted positive cell mean \\
Unskilled wage & \texttt{REGNOCALD} (\emph{no calificado}) & Weighted positive cell mean \\
Skill premium & \texttt{REGTESP} and \texttt{REGNOCALD} & $\log(\text{skilled})-\log(\text{unskilled})$ \\
\bottomrule
\end{tabular}
\begin{tablenotes}[flushleft]\footnotesize
\item \textit{Notes.} Field names refer to the harmonized wide ENIA files created from the annual raw
databases. The table reports the variables used in the main analysis; additional asset subclasses and
robustness variables are built in the same way.
\end{tablenotes}
\end{threeparttable}
\end{table}

\subsection{Earthquake intensity}
I measure earthquake exposure with the USGS ShakeMap for the February 27, 2010 Maule earthquake. I overlay
the ShakeMap MMI polygons on the GADM 4.1 administrative-region polygons, after reprojecting both layers to
a planar equal-area coordinate system so that the areas are computed correctly. Each region's intensity is
the area-weighted mean MMI:
\[ \text{MMI}_r = \frac{\sum_{g \in r} \text{MMI}_g \cdot \text{area}_{gr}}{\sum_{g \in r} \text{area}_{gr}}, \]
where $g$ indexes the ShakeMap polygons that intersect region $r$. For regions outside the felt area of the
earthquake, I assign a floor value of 2.0, which corresponds to weak or essentially unfelt shaking, so that
they serve as low-intensity controls rather than being dropped. To keep the region coding consistent across
the ENIA sample, I use the historical fifteen-region structure, so \~Nuble is still counted as part of
B\'io-B\'io.

\subsection{Detrending, weights, and inference}
The main threat to identification is that the hardest-hit central regions were already growing faster before
the earthquake. To deal with this, I estimate each region's linear trend using only the pre-earthquake years
(2003, 2004, 2005, 2006, 2007, 2008, and 2009), and I project this trend forward. I then run the event-study
and mechanism regressions on the deviations from this pre-period-projected path. I do not estimate the trend
on the whole sample, because that would let the post-earthquake reconstruction pulse affect the counterfactual
slope and pull the treatment effect toward zero.

I weight the regressions by cell effective size, clipping the cell weight at one for numerical stability, and
I cluster the standard errors by region. Because there are only fifteen region clusters, I read the
cluster-robust asymptotic standard errors cautiously, and for the headline investment pulse, the 2014
unskilled-wage response, and the equipment-accumulation check I also report restricted wild cluster bootstrap
$p$-values with Rademacher weights.

Finally, ENIA applies multiplicative confidentiality noise to the establishment records. This noise inflates
the variance of establishment-level measures, especially investment, but it does not create a systematic
difference-in-differences bias in expectation, as long as the perturbation is not correlated with earthquake
intensity conditional on the fixed effects. All of the tables and figures in the paper are produced by the
replication scripts, which go from the raw annual databases through extraction, harmonization, cell
construction, intensity merging, estimation, and table and figure generation.

\end{document}
